\subsection{\bt{Published memories on RelBench}{RelBench上的已发表记忆方法对照}}
\label{sec:relbench-eval}
\bi{\emph{Tasks and time separation.} We use RelBench's task definitions, label-generation code, and metrics~\cite{robinson2024relbench}: post-votes (MAE), user-badge and user-engagement (AUC) on rel-stack; site-success and study-adverse (MAE), and study-outcome (AUC) on rel-trial. The official-API study uses the first three cutoffs of our previously prepared five-cutoff grid: January, April, and July 2020 for rel-stack, and January 2017--2019 for rel-trial. These are earlier temporal evaluations, not the final official test split. Each task starts a fresh agent session at each cutoff; the next period's rows are appended between cutoffs. We schedule 90 sessions: five methods, two databases, three tasks, and three cutoffs.}
{\emph{任务与时间划分。}我们使用RelBench的任务定义、标签生成代码和指标~\cite{robinson2024relbench}：rel-stack上的post-votes（MAE）、user-badge与user-engagement（AUC）；rel-trial上的site-success与study-adverse（MAE）、study-outcome（AUC）。官方API实验采用既有五截点网格中的前三个：rel-stack为2020年1月、4月和7月，rel-trial为2017--2019年1月。这些是较早的时间评估，并非最终官方测试划分。每项任务在每个截点启动新的智能体会话，截点之间追加下一期的记录。共安排90个会话：五种方法、两个数据库、三个任务、三个截点。}

\bi{\emph{Development and baselines.} Before formal evaluation, we compare two predeclared MAVRA configurations using only rel-amazon/user-churn validation scores at two cutoffs. The selected policy first validates an available task program, retrieves feature definitions as needed, and skips repeated extraction only for identical SQL already published and freshly audited. Its configuration is frozen before formal test scoring. The two development candidates consume 0.669\,M tokens in total; this one-time tuning cost is reported separately from the 90 formal sessions and is included in the study's API spending. The comparators are the same agent without memory, ACE~\cite{zhang2026ace}, AWM~\cite{wang2025awm}, and G-Memory~\cite{zhang2025gmemory}. ACE maintains a playbook per task using its official reflector and curator; AWM induces task workflows from successful sessions using the official prompt; G-Memory uses the official implementation with a database-wide shared memory. We adapt their published code and prompts to the common SQL tool interface and give all methods the same validation feedback. Reproduction details, revisions, and adaptations are in the artifact. Earlier gateway-based runs, including DeepPrep and FeatPilot comparisons, remain archived separately and are not pooled with this study.}
{\emph{开发与基线。}正式评估前，我们仅用rel-amazon/user-churn两个截点的验证分数比较两种预先确定的MAVRA配置。选中的策略先验证可用的任务方案，按需检索特征定义，并仅对已发布且本次重新通过审计的逐字相同SQL跳过重复提取。配置在正式测试评分前冻结。两个开发候选总计使用0.669\,M token；该一次性调优成本与90个正式会话分列，并计入本轮API支出。对照为同一智能体的无记忆版本，以及ACE~\cite{zhang2026ace}、AWM~\cite{wang2025awm}和G-Memory~\cite{zhang2025gmemory}。ACE用官方反思者与整理者维护每任务手册；AWM用官方提示词从成功会话归纳任务工作流；G-Memory使用官方实现，在同一数据库的智能体之间共享记忆。我们将其公开代码与提示词适配到共同的SQL工具接口，各方法获得相同类型的验证反馈。复现版本与适配细节见工件。较早的网关运行及其中的DeepPrep、FeatPilot对照单独归档，不与本轮合并。}

\bi{\emph{Common execution and accounting.} All methods use the official DeepSeek endpoint with the \code{deepseek-flash} alias, thinking enabled, and low reasoning effort; the provider documents this alias as \llmmodel{} on the study date. Agents share schema, join, read-only SQL, validation, and submission tools. Feature SQL returns one row per seed using only data no later than its time. Every validation fits a fixed LightGBM~\cite{ke2017lightgbm} on 50,000 sampled training seeds and performs a 5,000-seed temporal-leakage audit; only explicitly clean audits can supply a valid solution. The limits are 24 planner calls, five validations, two non-improving validations, and 600 seconds per feature SQL statement. From cutoff 2, the previous solution refit on current data supplies the validation target; after reaching it, one further improvement is allowed within three calls. Test labels do not enter agent feedback or memory updates, and final test scoring follows the sequence. Failed solutions receive the entity-table reference and retain all incurred costs. Total tokens include memory retrieval, extraction, and updates; reuse totals cover cutoffs 2--3. Runs share a server, so observed time is not an isolated speedup measure.}
{\emph{统一执行与计量。}所有方法均通过DeepSeek官方端点调用\code{deepseek-flash}，启用思考模式、低推理强度；提供方在实验日期将该别名记为\llmmodel{}。智能体共用看结构、连接、只读SQL、验证与提交工具。特征SQL为每个种子返回一行，只能使用不晚于种子时间的数据。每次验证在5万个抽样训练种子上拟合固定LightGBM~\cite{ke2017lightgbm}，并进行5000种子的时间泄漏审计；只有明确clean的审计结果可以提供有效方案。上限为24次规划模型调用、5次验证、连续2次不改进，以及每条特征SQL语句600秒。从截点2起，用上一方案在当前数据上的重新拟合分数作为验证目标；达到目标后，可在3次调用内再改进一次。测试标签不进入智能体反馈或记忆更新，最终测试评分在序列结束后执行。失败方案使用实体表参照分数并保留全部已发生成本。总token计入记忆检索、提取和更新，复用代价覆盖截点2--3。各组共享试验机，因此观测耗时不作为隔离加速比。}

\input{tables/official-results}
\begin{figure}[t]
\centering
\ifbilingual
\includegraphics[width=\linewidth]{figures/official-cumulative-zh}
\else
\includegraphics[width=\linewidth]{figures/official-cumulative}
\fi
\caption{\bt{Cumulative tokens through each retraining cutoff, including all three tasks, unsuccessful attempts, and memory retrieval and updates. Cutoff 1 includes initial construction; cutoffs 2--3 permit reuse. Each line is one sequence, not a mean over independent repetitions.}{截至每个重训截点的累计token，包含三个任务、失败尝试以及记忆检索与更新。截点1计入初始构建，截点2--3允许复用。每条线为一条序列，并非多次独立重复的均值。}}
\label{fig:official-cumulative}
\Description{Two vertically stacked line charts show cumulative tokens in millions at cutoffs 1, 2 and 3 for all five methods on rel-stack and rel-trial. Distinct colors and markers identify methods; both y-axes start at zero.}
\end{figure}

\bi{\emph{Quality and amortization.} Table~\ref{tab:official-results} reports every task's mean test quality, valid-session counts, and costs; Figure~\ref{fig:official-cumulative} charges memory construction from the first cutoff. Relative to no memory, \system\ reduces total tokens by \OfStackSaveAll\% on rel-stack and \OfTrialSaveAll\% on rel-trial, and reuse-stage tokens by \OfStackSaveReuse\% and \OfTrialSaveReuse\%. Across all cutoffs, \OfFullQualityCount/6 tasks meet the predeclared quality margins, $\mathrm{AUC}\geq\mathrm{AUC}_0-0.01$ or $\mathrm{MAE}\leq1.02\,\mathrm{MAE}_0$. The exception is study-adverse, whose full-sequence MAE is \OfAdverseRegression\% higher; its first-cutoff attempt has no valid solution. All six tasks meet the same margins over reuse cutoffs, while those first-cutoff failures and their costs remain in the full-sequence results. These are descriptive results from one sequence per method: cutoffs are dependent observations, not independent repetitions.}
{\emph{质量与成本分摊。}表~\ref{tab:official-results}报告每项任务的平均测试质量、有效会话数和代价；图~\ref{fig:official-cumulative}从首截点起计入记忆构建。相对无记忆，\system\ 在rel-stack与rel-trial上的总token分别减少\OfStackSaveAll\%与\OfTrialSaveAll\%，复用阶段分别减少\OfStackSaveReuse\%与\OfTrialSaveReuse\%。全部截点上，\OfFullQualityCount/6项任务满足预设质量门槛，即$\mathrm{AUC}\geq\mathrm{AUC}_0-0.01$或$\mathrm{MAE}\leq1.02\,\mathrm{MAE}_0$。例外是study-adverse，全序列MAE高\OfAdverseRegression\%，其首截点尝试没有得到有效方案。复用截点上六项任务均满足相同门槛；首截点失败及其代价仍计入全序列结果。每种方法仅有一条序列，这些结果是描述性比较，截点彼此依赖，并非独立重复。}


\bi{\emph{Comparison with published memories.} Relative to ACE, AWM, and G-Memory, respectively, \system\ uses 45.7\%, 49.9\%, and 65.1\% fewer total tokens on rel-stack, and 36.7\%, 33.1\%, and 49.5\% fewer on rel-trial. Quality remains task-dependent: G-Memory has lower post-votes MAE and higher user-badge AUC on rel-stack, and lower study-adverse MAE on rel-trial. AWM retains seven valid rel-trial sessions versus six for \system. Table~\ref{tab:official-results} therefore reports quality and failures alongside the cost reductions.}
{\emph{与已发表记忆方法比较。}相对ACE、AWM和G-Memory，\system\ 在rel-stack上的总token分别减少45.7\%、49.9\%和65.1\%，在rel-trial上分别减少36.7\%、33.1\%和49.5\%。质量仍依任务而异：G-Memory在rel-stack上的post-votes MAE更低、user-badge AUC更高，在rel-trial上的study-adverse MAE更低。AWM在rel-trial保留七个有效会话，\system\ 为六个。因此，表~\ref{tab:official-results}将质量与失败次数同成本节省并列报告。}
